\chapter{Materiales y Métodos}

\section{Diseño de la investigación}

La presente investigación se inscribe en un enfoque cuantitativo, con un diseño no experimental, observacional, retrospectivo, longitudinal y explicativo. Es de carácter observacional y no experimental debido a que las variables climáticas se analizan de forma pasiva, tal como ocurrieron de manera natural en el entorno, sin manipulación deliberada. Posee un enfoque retrospectivo, puesto que se examinan registros históricos de un período extendido de 24 años (2000--2023). Asimismo, su diseño es longitudinal y explicativo mediante el uso de una estructura de datos de panel, orientada a estimar parámetros de tendencia de largo plazo para explicar y contrastar formalmente las hipótesis sobre las dinámicas de la sequía y el proceso de aridificación en Paraguay.

\section{Fuentes de datos y variables}

Los datos primarios corresponden a los registros meteorológicos diarios provistos por la Dirección de Meteorología e Hidrología (DMH) de la Dirección Nacional de Aeronáutica Civil (DINAC) del Paraguay. 

\subsection{Procedimiento de extracción y validación de Datos}

La serie temporal diaria analizada comprende el período desde el 1 de enero de 2000 hasta el 31 de diciembre de 2023. La fecha de extracción y recepción oficial de la base de datos fue el 9 de mayo de 2024 (Expediente N.º 213.042/2024; ver Apéndice \ref{apendice:fuente_datos}), tras una solicitud formal gestionada a través de la Facultad de Ciencias Exactas y Naturales (FACEN) de la Universidad Nacional de Asunción.

De forma complementaria, se recopilaron los metadatos de las estaciones terrestres, incluyendo su nombre oficial, identificador (ID Estación), coordenadas geográficas de posicionamiento (latitud y longitud decimal) y su macroregión geográfica correspondiente (Región Oriental o Región Occidental).

\subsection{Definición de variables climáticas}

El análisis de datos de panel se centra de manera exclusiva en las variables de precipitación diaria (lluvia) y temperatura media diaria como factores determinantes del balance hídrico. Las variables críticas recopiladas con frecuencia diaria para cada estación meteorológica fueron:
\begin{itemize}
    \item \textit{Precipitación diaria ($P_t$):} medida en milímetros ($mm$).
    \item \textit{Temperatura media diaria ($T_{\text{med}, t}$):} medida en grados Celsius ($^{\circ}\text{C}$).
    \item \textit{Temperatura máxima diaria ($T_{\text{máx}, t}$):} usada como variable de apoyo para la imputación térmica ($^{\circ}\text{C}$).
    \item \textit{Temperatura mínima diaria ($T_{\text{mín}, t}$):} usada como variable de apoyo para la imputación térmica ($^{\circ}\text{C}$).
\end{itemize}

\section{Procedimiento, depuración e imputación de datos}

Se examinó la cobertura histórica de cada una de las 19 estaciones terrestres respecto al rango teórico esperado de 8,766 días (24 años). El análisis de completitud determinó que la estación de Paraguarí (PARA) presentaba un truncamiento temporal crítico: sus registros inician únicamente a partir del 1 de abril de 2006, registrando una pérdida masiva superior al 26\% de la serie temporal bajo estudio (2,282 días faltantes). 

Con el fin de evitar un sesgo por truncamiento sistemático de la serie temporal nacional, se decidió excluir de manera definitiva la estación de Paraguarí (PARA). Las 18 estaciones restantes exhibieron tasas de completitud de fechas superiores al 97.9\%, conformando la muestra final del estudio. Sobre esta red de 18 estaciones se construyó un panel balanceado inicial de 157,788 registros diarios (18 estaciones $\times$ 8,766 días). Las filas correspondientes a fechas faltantes originalmente en cada estación fueron insertadas e inicializadas con valores nulos (\textit{NaN}).

\subsection{Algoritmo de imputación de temperatura por capas}

Para asegurar la continuidad temporal de la temperatura media diaria ($T {\text{med}}$) y evitar la pérdida de observaciones mensuales, se diseñó e implementó un algoritmo de imputación estructurado en tres capas jerárquicas:

\begin{enumerate}
    \item \textit{Capa 1: Imputación por relación térmica.} Para aquellos días donde se registró la ausencia de la temperatura media ($T {\text{med}}$) pero se disponía de los registros válidos de temperatura máxima ($T {\text{máx}}$) y mínima ($T {\text{mín}}$), se estimó el valor faltante como el promedio aritmético de las temperaturas extremas diarias:
    \begin{equation}
        T_{\text{med}, t} = \frac{T_{\text{mín}, t} + T_{\text{máx}, t}}{2}
    \end{equation}
    Este procedimiento permitió recuperar con precisión física un total de 342 días del panel.
    \item \textit{Capa 2: Interpolación lineal para vacíos cortos ($\le 3$ días).} Para los registros nulos persistentes de temperatura media, se analizó la longitud de los vacíos consecutivos. Para brechas de duración corta, definidas estrictamente entre 1 y 3 días consecutivos (que representan el 37.1\% del total de datos térmicos faltantes en la red), se aplicó una interpolación lineal temporal basada en la inercia térmica de los días adyacentes válidos. Esta capa recuperó un total de 655 días de registro.
    \item \textit{Capa 3: Imputación espacial por estación vecina ($\geq 4$ días).} Para vacíos de información térmica igual o mayor a cuatro días consecutivos, la interpolación temporal directa pierde validez física al ignorar la variabilidad sinóptica de mediano plazo. En consecuencia, se aplicó un método de transferencia espacial basado en la cercanía geográfica. 

Utilizando las coordenadas geográficas de las 18 estaciones meteorológicas seleccionadas, se calculó una matriz de distancias mutuas (ortodrómicas) que permitió identificar de forma unívoca a la estación vecina más cercana para cada punto de la red (por ejemplo, la menor distancia registrada correspondió al par Coronel Oviedo-Villarrica con solo $25.75\text{ km}$, mientras que la mayor distancia de vecindad se registró en el Chaco Profundo entre Mariscal Estigarribia y Pozo Colorado con $248.84\text{ km}$). 

Para cada fecha que presentaba vacíos térmicos prolongados en la estación objetiva, se transfirió directamente el registro diario de temperatura media depurada ($T_{\text{med\_imputada}}$) correspondiente a su vecina más cercana en ese mismo día. Este procedimiento de interpolación espacial permitió recuperar con precisión climatológica un total de 1,114 observaciones diarias. Al término de esta tercera capa jerárquica de imputación, se logró una tasa de completitud global en las series de temperatura media diaria del 99.99\%, registrando únicamente 9 datos nulos residuales en toda la base de datos de panel maestro para las 18 estaciones a lo largo de los 24 años de estudio.
\end{enumerate}

\subsection{Criterio de no Imputación para la variable precipitación}

A diferencia de la temperatura, la precipitación diaria es una variable altamente asimétrica, intermitente y no continua. Debido a que las fallas instrumentales y los vacíos de transmisión de la red meteorológica ocurren por factores técnicos o administrativos totalmente independientes del estado del tiempo, no existe relación causal entre la omisión del dato y el evento físico de lluvia. Bajo esta premisa, como lo afirma \citeA{wooldridge2010}, el sesgo no depende de la cantidad o el patrón de los datos faltantes, sino de la naturaleza del mecanismo de selección que genera esos nulos. Si la falta de datos es aleatoria o técnica, el estimador de efectos fijos manejará correctamente el desbalance del panel sin introducir sesgos de truncamiento temporal.

Consecuentemente, se adoptó la decisión metodológica de no aplicar ningún método de imputación sobre la variable precipitación. Forzar estimaciones (ya sean interpolaciones lineales o vecindades) sobre la lluvia introduciría sesgos artificiales severos, alterando la frecuencia real de días lluviosos y la magnitud acumulada de los eventos.

\section{Agregación de datos y creación de la base mensual}

Dado que el balance hídrico y el cálculo del SPEI se analizan a nivel mensual, el panel diario depurado de 18 estaciones fue agregado a una frecuencia mensual de 288 meses consecutivos por estación (5,184 observaciones mensuales potenciales). 

Para garantizar la fiabilidad física de los agregados mensuales, se implementó la regla de exclusión recomendada por la Organización Meteorológica Mundial \cite{wmo2017}:
\begin{itemize}
    \item \textit{lluvia\_total:} suma mensual de la precipitación diaria. Se define como nula (\textit{NaN}) si el mes correspondiente registra menos de 25 días con observaciones de lluvia completas.
    \item \textit{temp\_media:} promedio mensual de la temperatura media diaria imputada. Se define como nula si el mes registra menos de 25 días con datos térmicos válidos.
\end{itemize}

Este procedimiento determinó un panel mensual final caracterizado por una alta completitud: la serie de temperatura media mensual reportó un único dato nulo ($0.01\%$), mientras que la serie de precipitación mensual acumulada presentó únicamente 82 registros nulos ($1.58\%$) de un total de 5,184 observaciones.

\section{Construcción de la variable dependiente: el índice SPEI-6}

Para el cálculo del Índice de Precipitación-Evapotranspiración Estandarizado a escala de 6 meses (SPEI-6) sobre las 18 estaciones finales, se ejecutó secuencialmente la siguiente formulación matemática:

\begin{enumerate}
    \item \textit{Cálculo de la ETP mensual (Thornthwaite):} Estimación de la demanda evaporativa potencial mediante el método empírico de \citeA{thornthwaite_approach_1948}, utilizando la temperatura media mensual ($T_{\text{avg}}$) y la latitud geográfica de cada estación para derivar el exponente empírico $a$ y el factor de corrección latitudinal mensual $K$.
    \item \textit{Cálculo del balance hídrico mensual ($D_i$):} Definido como la diferencia simple entre la precipitación total mensual ($P_i$) y la ETP ajustada ($ETP_i$) para cada mes $i$: $D_i = P_i - ETP_i$.
    \item \textit{Acumulación multiescalar ($D^k_i$):} Para capturar los efectos de mediano plazo se acumuló el balance hídrico mensual en una ventana temporal de $k=6$ meses de forma consecutiva \cite{mckee1993}.
    \item \textit{Ajuste paramétrico Pearson Tipo III:} Los valores acumulados de balance hídrico ($D^6_i$) fueron ajustados a una distribución teórica Pearson Tipo III (Gamma de tres parámetros) para estimar la probabilidad acumulada $P = F(x)$, lo cual resulta óptimo para tolerar asimetrías y valores negativos de balance.
    \item \textit{Estandarización Normal estándar:} La probabilidad acumulada $P$ se transformó en anomalías estandarizadas normales ($Z$-scores) mediante la aproximación racional analítica de \citeA{abramowitz_handbook_1965}. El valor final resultante constituye la variable dependiente de panel, SPEI-6.
\end{enumerate}

\section{Regionalización geográfica mediante agrupamiento no supervisado}

Para analizar la distribución espacial de la sequía y evitar límites políticos arbitrarios, se aplicó un algoritmo de agrupamiento jerárquico aglomerativo (HAC) sobre las series de SPEI-6 de las 18 estaciones. Se implementó la distancia de disimilitud euclídea y el criterio de enlace de varianza mínima de Ward. El número óptimo de subregiones climáticas ($k^*$) se determinó mediante la maximización del Coeficiente de Silueta Promedio de Rousseeuw \cite{rousseeuw1987}.

\section{Especificación del modelo de panel y test de Hausman}\label{sec:modelo_estadistico}

Para cuantificar las tendencias temporales del balance hídrico se ajustó el modelo de regresión lineal de panel especificado en la ecuación \ref{eq:modelo_completo} y la selección entre la estimación por efectos fijos (FE) y efectos aleatorios (RE) se fundamentó en el test de especificación de \citeA{hausman1978}, evaluando formalmente la consistencia del estimador RE frente a la ortogonalidad de los efectos individuales no observados de las 18 estaciones meteorológicas. Para esto se utilizó la librería \textit{linearmodels} de Python.

La especificación completa del modelo estimado, incluyendo todas las variables e interacciones, fue la siguiente:
\sloppy
\begin{equation}
\begin{split}
    \text{SPEI-6}_{it} = & \beta_0 + \beta_1 \text{tiempo}_t + \sum_{s=1}^{3} \gamma_s \text{estacion-climatica}_{s,t} + \delta_1 \text{region-geografica}_i \\
    & + \sum_{s=1}^{3} \theta_s (\text{tiempo}_t \times \text{estacion-climatica}_{s,t}) + \phi_1 (\text{tiempo}_t \times \text{region-geografica}_i) \\
    & + \sum_{s=1}^{3} \psi_s (\text{estacion-climatica}_{s,t} \times \text{region-geografica}_i) \\
    & + \sum_{s=1}^{3} \omega_s (\text{tiempo}_t \times \text{estacion-climatica}_{s,t} \times \text{region-geografica}_i) + u_{it}
\end{split}
\label{eq:modelo_completo}
\end{equation}
donde $u_{it} = \alpha_i + \varepsilon_{it}$ es el término de error compuesto, y los términos $\text{estacion-climatica}_{s,t}$ y $\text{region-geografica}_i$ son variables dummys para las categorías estacionales climáticas (invierno, primavera, verano, otoño) y regionales geográficas (clusters), respectivamente.

\subsection{Corrección de la dependencia espacial: el estimador de Driscoll y Kraay}

Para controlar y corregir de manera simultánea la presencia de autocorrelación serial temporal y de correlación espacial o de sección cruzada identificada en la red de monitoreo, se prescinde de los errores estándar de clúster convencionales y se adopta el estimador de matriz de covarianza robusto desarrollado por \cite{driscoll1998}. 

El método de Driscoll-Kraay es un estimador no paramétrico de matriz de covarianza consistente ante la presencia de heterocedasticidad, autocorrelación y cualquier forma de dependencia espacial o de sección cruzada. Su ventaja metodológica radica en que estima la estructura de covarianza de las medias temporales de los momentos transversales mediante un estimador tipo Newey-West con un kernel espectral de Bartlett, asegurando la consistencia de los errores estándar y garantizando intervalos de confianza robustos para la inferencia de las tendencias dinámicas de sequía a nivel nacional \cite{wooldridge2010}.

\section{Marco de contraste SPEI vs. SPI para el aislamiento del efecto térmico}

Para dar cumplimiento a nuestro objetivo específico y contrastar de manera formal la hipótesis 3 (H3), se diseñó un marco experimental comparativo. El procedimiento consiste en estimar de forma paralela la ecuación del modelo de regresión de panel de efectos aleatorios (Ecuación \ref{eq:modelo_completo}) sobre dos variables dependientes distintas:
\begin{align}
    \text{SPEI-6}_{it} &= \mathbf{x}_{it}'\boldsymbol{\beta}_{SPEI} + u_{it} \label{eq:panel_spei} \\
    \text{SPI-6}_{it} &= \mathbf{x}_{it}'\boldsymbol{\beta}_{SPI} + v_{it} \label{eq:panel_spi}
\end{align}

Donde $\boldsymbol{\beta}_{SPEI}$ y $\boldsymbol{\beta}_{SPI}$ representan los vectores de parámetros estimados para los modelos de balance hídrico y precipitación univariada, respectivamente. El efecto aislado de la ETP sobre la dinámica temporal del balance hídrico se define para cada estación del año y región como la diferencia entre sus respectivas tendencias netas mensuales calculadas ($\Delta \beta_{neto} = \beta_{neto, SPEI} - \beta_{neto, SPI}$). 

La significancia estadística del efecto del aumento de la demanda evaporativa impulsada por el aumento de las temperaturas (calentamiento global) ($H_0: \beta_{neto, SPEI} - \beta_{neto, SPI} = 0$) se evalúa mediante la estimación de intervalos de confianza del 95\% y pruebas de Wald para diferencias de coeficientes de pendiente comunes entre ambas regresiones de panel, aislando así el efecto puro de la temperatura por encima del régimen de precipitaciones.
